\QHeading{第1题　在\(\bigcirc\)和\(\square\)里分别填写相应的运算符号和数。}
\par\noindent


\PlainItem{\(17\times48=48\times\){\color{answerred}17}\(50\times76\times2=(\){\color{answerred}50}\(\times\){\color{answerred}2}\()\times76\)}
\par

\PlainItem{\(34\times\){\color{answerred}19}\(=19\times\){\color{answerred}34}\(39\times5\times40=39\times(\){\color{answerred}5}{\color{answerred}\(\times\)}{\color{answerred}40}\()\)}
\PlainItem{\(25\times5\times4\times2=(\){\color{answerred}25}\(\times\){\color{answerred}4}\()\times(\){\color{answerred}5}\(\times\){\color{answerred}2}\()\)}
\par\noindent

\QHeading{第2题　用简便方法计算下面各题。}

\CalcBlock[方法一：简便计算]{1}{\begin{aligned}[t]&\phantom{={}}\color{black}25\times66\times4\\&\color{answerred}=(25\times4)\times66\\&\color{answerred}=100\times66\\&\color{answerred}=6600\end{aligned}}
\par
\MethodBlock{方法二：结合因数，先凑整}{用乘法交换律、结合律，先算容易凑整的两个因数。}{\begin{aligned}[t]&\phantom{={}}25\times 66\times 4\\
&=(25\times 66)\times 4\\
&=1650\times 4\\
&=6600\end{aligned}}


\CalcBlock[方法一：简便计算]{2}{\begin{aligned}[t]&\phantom{={}}\color{black}24\times25\times8\\&\color{answerred}=24\times(25\times8)\\&\color{answerred}=24\times200\\&\color{answerred}=4800\end{aligned}}
\MethodBlock{方法二：结合因数，先凑整}{用乘法交换律、结合律，先算容易凑整的两个因数。}{\begin{aligned}[t]&\phantom{={}}24\times 25\times 8\\
&=(24\times 25)\times 8\\
&=600\times 8\\
&=4800\end{aligned}}

\CalcBlock[方法一：简便计算]{3}{\begin{aligned}[t]&\phantom{={}}\color{black}4\times173\times5\\&\color{answerred}=(4\times5)\times173\\&\color{answerred}=20\times173\\&\color{answerred}=3460\end{aligned}}
\MethodBlock{方法二：先处理括号或乘除，再凑整后补偿}{把接近整十、整百的因数凑整，再补回或减去多算的部分。}{\begin{aligned}[t]&\phantom{={}}4\times 173\times 5\\
&=692\times 5\\
&=5\times (690+2)\\
&=5\times 690+5\times 2\\
&=3450+10\\
&=3460\end{aligned}}

\CalcBlock[方法一：简便计算]{4}{\begin{aligned}[t]&\phantom{={}}\color{black}4\times36\times25\\&\color{answerred}=(4\times25)\times36\\&\color{answerred}=100\times36\\&\color{answerred}=3600\end{aligned}}
\MethodBlock{方法二：结合因数，先凑整}{用乘法交换律、结合律，先算容易凑整的两个因数。}{\begin{aligned}[t]&\phantom{={}}4\times 36\times 25\\
&=(36\times 25)\times 4\\
&=900\times 4\\
&=3600\end{aligned}}

\MethodBlock{方法三：先处理括号或乘除，再把特殊因数改成整百或整千}{先把25、125等改写成整百或整千除以一个数。}{\begin{aligned}[t]&\phantom{={}}4\times 36\times 25\\
&=144\times 25\\
&=(100\div 4)\times 144\\
&=100\times (144\div 4)\\
&=100\times 36\\
&=3600\end{aligned}}

\CalcBlock[方法一：简便计算]{5}{\begin{aligned}[t]&\phantom{={}}\color{black}7\times5\times12\\&\color{answerred}=7\times(5\times12)\\&\color{answerred}=7\times60\\&\color{answerred}=420\end{aligned}}
\MethodBlock{方法二：先处理括号或乘除，再凑整后补偿}{把接近整十、整百的因数凑整，再补回或减去多算的部分。}{\begin{aligned}[t]&\phantom{={}}7\times 5\times 12\\
&=35\times 12\\
&=35\times (10+2)\\
&=35\times 10+35\times 2\\
&=350+70\\
&=420\end{aligned}}

\CalcBlock[方法一：简便计算]{6}{\begin{aligned}[t]&\phantom{={}}\color{black}125\times19\times8\\&\color{answerred}=(125\times8)\times19\\&\color{answerred}=1000\times19\\&\color{answerred}=19000\end{aligned}}
\MethodBlock{方法二：先处理括号或乘除，再凑整后补偿}{把接近整十、整百的因数凑整，再补回或减去多算的部分。}{\begin{aligned}[t]&\phantom{={}}125\times 19\times 8\\
&=2375\times 8\\
&=2375\times (10-2)\\
&=2375\times 10-2375\times 2\\
&=23750-4750\\
&=19000\end{aligned}}

\CalcBlock[方法一：简便计算]{7}{\begin{aligned}[t]&\phantom{={}}\color{black}125\times(23\times8)\\&\color{answerred}=(125\times8)\times23\\&\color{answerred}=1000\times23\\&\color{answerred}=23000\end{aligned}}
\MethodBlock{方法二：先处理括号或乘除，再把特殊因数改成整百或整千}{先把25、125等改写成整百或整千除以一个数。}{\begin{aligned}[t]&\phantom{={}}125\times (23\times 8)\\
&=125\times 184\\
&=(1000\div 8)\times 184\\
&=1000\times (184\div 8)\\
&=1000\times 23\\
&=23000\end{aligned}}

\CalcBlock[方法一：简便计算]{8}{\begin{aligned}[t]&\phantom{={}}\color{black}25\times(82\times4)\\&\color{answerred}=(25\times4)\times82\\&\color{answerred}=100\times82\\&\color{answerred}=8200\end{aligned}}
\MethodBlock{方法二：结合因数，先凑整}{用乘法交换律、结合律，先算容易凑整的两个因数。}{\begin{aligned}[t]&\phantom{={}}25\times (82\times 4)\\
&=(25\times 82)\times 4\\
&=2050\times 4\\
&=8200\end{aligned}}

\MethodBlock{方法三：先处理括号或乘除，再把特殊因数改成整百或整千}{先把25、125等改写成整百或整千除以一个数。}{\begin{aligned}[t]&\phantom{={}}25\times (82\times 4)\\
&=25\times 328\\
&=(100\div 4)\times 328\\
&=100\times (328\div 4)\\
&=100\times 82\\
&=8200\end{aligned}}

\MethodBlock{方法四：先处理括号或乘除，再凑整后补偿}{把接近整十、整百的因数凑整，再补回或减去多算的部分。}{\begin{aligned}[t]&\phantom{={}}25\times (82\times 4)\\
&=25\times 328\\
&=25\times (330-2)\\
&=25\times 330-25\times 2\\
&=8250-50\\
&=8200\end{aligned}}

\CalcBlock[方法一：简便计算]{9}{\begin{aligned}[t]&\phantom{={}}\color{black}125\times48\\&\color{answerred}=125\times(8\times6)\\&\color{answerred}=(125\times8)\times6\\&\color{answerred}=1000\times6\\&\color{answerred}=6000\end{aligned}}
\MethodBlock{方法二：把特殊因数改成整百或整千}{先把25、125等改写成整百或整千除以一个数。}{\begin{aligned}[t]&\phantom{={}}125\times 48\\
&=(1000\div 8)\times 48\\
&=1000\times (48\div 8)\\
&=1000\times 6\\
&=6000\end{aligned}}

\MethodBlock{方法三：凑整后补偿}{把接近整十、整百的因数凑整，再补回或减去多算的部分。}{\begin{aligned}[t]&\phantom{={}}125\times 48\\
&=125\times (50-2)\\
&=125\times 50-125\times 2\\
&=6250-250\\
&=6000\end{aligned}}

\MethodBlock{方法四：按数位拆分计算}{把一个因数拆成整十、整百数与余下部分之和。}{\begin{aligned}[t]&\phantom{={}}125\times 48\\
&=125\times (40+8)\\
&=125\times 40+125\times 8\\
&=5000+1000\\
&=6000\end{aligned}}

\par\noindent

\QHeading{第3题　解决问题。}

\Problem{学校图书室有17个书架，每个书架有5层，平均每层有20本书。学校图书室一共有多少本书？}
\par


\MethodTitle{方法一}
\Answer{\Formula{\begin{aligned}[t]&\phantom{={}}20\times5\times17\\&=100\times17\\&=1700\text{（本）}\end{aligned}}\par 答：学校图书室一共有1700本书。}

\MethodTitle{方法二：先求全部书架的总层数}
\Answer{共有17个书架，每个5层，再按每层20本计算。}
\Answer{\Formula{\begin{aligned}[t]&\phantom{={}}20\times (5\times 17)\\
&=20\times 85\\
&=1700\text{（本）}\end{aligned}}}
\Answer{答：图书室一共有1700本书。}
